\chapter{The derivative}

\section{Mean value theorem}
\label{sec:mvt}

\begin{defn}
A function $f$ has a \emph{relative maximum} at $c$ when some interval
around $c$ has no point $x$ with $f(x) > f(c)$.
\end{defn}

\begin{lemma} \label{lemma:relmaxdzero}
If $f$ is differentiable at a relative maximum $c$, then $f'(c) = 0$.
\end{lemma}

\begin{proof}
The difference quotients on the two sides of $c$ have opposite signs.
\end{proof}

\begin{thm}[Rolle]
\label{thm:rolle}
If $f$ is continuous on $[a,b]$, differentiable on $(a,b)$ and $f(a) = f(b)$,
then $f'(c) = 0$ for some $c$ in $(a,b)$.
\end{thm}

\begin{proof}
The function has an extreme value inside the interval, and
\lemmaref{lemma:relmaxdzero} applies there:
\begin{equation}
\label{eq:rolle-zero}
f'(c) = 0 .
\end{equation}
\end{proof}

\begin{example}
The function $x^2$ on $[-1,1]$ shows the theorem with $c = 0$.
\end{example}

\subsection{Exercises}

\begin{exercise}
Find a function where the hypothesis of \ref{thm:rolle} fails and the
conclusion fails too.
\end{exercise}

\begin{exercise}
Prove \hyperref[thm:rolle]{Rolle's theorem} again with a different extreme value.
\end{exercise}
